\section{Explanation in Artificial Intelligence}
\label{sec:explanation}

In this section, we define what an explanation is in the context of \gls{xai} and how it can be obtained.
We first introduce the concept of explanation in general, along with the explainability method used to produce it.
We then describe how this concept applies to \gls{xai}.
Finally, we further detail the properties that characterize an explainability method in the context of \gls{xai}.

\subsection{Explanation}
Establishing what is meant by an explanation is challenging due to the fundamentally multidisciplinary nature of the concept.
When we refer to explanation, we are inherently referring to an explanation provided to a human.
As soon as humans are involved, the problem is no longer purely computational or mathematical.
It involves several fields of research, such as philosophy, psychology, and sociology.
So far, these disciplines have not converged on a unified definition of explanation that can be directly applied to \gls{xai} \cite{miller2019explanation, bellucci2021towards}.

For this reason, we propose the following definition of explanation.
An explanation involves two entities: the object (Definition~\ref{def:object}) being explained and the receiver of the explanation.
The receiver is a human who aims to understand how the object operates and will be referred to as the observer throughout this manuscript (Definition~\ref{def:observer}).
The observer seeks to understand the object and, in doing so, builds a mental model of its functioning.
An explanation is therefore what allows the observer to refine and correct their mental model of how the object operates (Definition~\ref{def:explanation}).
This refinement is obtained through a transformation process that produces an explanation from the object being explained.
We refer to this process as an explainability method (Definition~\ref{def:explainability_method}), detailed further in Section~\ref{subsec:explainability_method}.

For an observer to understand an explanation, the information it contains must be interpretable (Definition~\ref{def:interpretable}).
Information is considered interpretable when it can be directly understood by a human without requiring any additional transformation process.

\begin{definition}{Object}{object}
The object of an explanation is the system the observer wants to understand.
\end{definition}

\begin{definition}{Observer}{observer}
The observer is the human who aims to understand an object.
The observer possesses a mental model of the object being explained.
This mental model is an internal representation of how the observer believes the system operates.
This representation may be incomplete or inaccurate.
\end{definition}

\begin{definition}{Interpretable}{interpretable}
Information is interpretable when it can be directly understood by an observer without requiring any additional transformation process.
\end{definition}

\begin{definition}{Explanation}{explanation}
An explanation is any interpretable information that enables an observer to refine or correct their mental model of the object of the explanation.
\end{definition}

\begin{definition}{Explainability Method}{explainability_method}
An explainability method is a transformation process that produces an explanation from the object being explained.
\end{definition}

This definition of explanation provides the foundation for the approach to explainability adopted throughout this manuscript.
Its generality allow it to encompass the different explainability methods considered in the following sections.

Moreover, this perspective avoids a binary interpretation of explainability, in which a system is considered either explainable or non-explainable.
The notion of a non-explainable system is difficult to justify from a scientific standpoint.
Scientific inquiry is driven by the assumption that observed phenomena can be investigated through the development of increasingly accurate analytical tools.
Declaring a system to be inherently impossible to explain would therefore conflict with this methodological principle.
This approach aligns with the philosophical Principle of Sufficient Reason, which states that every phenomenon rests upon an explainable basis \cite{sep-sufficient-reason}.

This quantitative perspective allows us to express an intuitive idea: not all explanations provide the same level of value \cite{doshivelez2017rigorous}.
The effectiveness of an explanation depends on its ability to improve the observer's understanding of the system.
For example, describing a marginal aspect of a system is different from correcting an observer's misunderstanding of a key mechanism.

\begin{figure}[t]
    \centering

    % --- Layout parameters --------------------------------------------
    % Horizontal distance between the centers of two consecutive blocks.
    \pgfmathsetmacro{\blockDistance}{6.6}
    % Width/height of each block (keep in sync with the node styles below).
    \pgfmathsetmacro{\blockSize}{3.6}
    % Fixed text width inside each block, so that longer labels wrap onto
    % extra lines (growing the block downward) instead of growing wider
    % and colliding with the neighboring block.
    \pgfmathsetmacro{\blockTextWidth}{3.1}
    % Text width used for the labels placed above/below the arrows.
    \pgfmathsetmacro{\arrowTextWidth}{2.6}

    % Vertical offset of the legend row below the main diagram.
    \pgfmathsetmacro{\legendYShift}{-3.6}
    % Horizontal offset of the start of the legend row (position of its first icon).
    \pgfmathsetmacro{\legendXShift}{-\blockSize/2 + 0}
    % Horizontal gap between two consecutive legend entries.
    \pgfmathsetmacro{\legendGap}{1.5}
    % Horizontal gap between a legend icon and its own label.
    \pgfmathsetmacro{\legendIconTextGap}{0.3}
    % --------------------------------------------------------------------

\resizebox{\textwidth}{!}{
\begin{tikzpicture}[]
% Nodes
\node[draw, rectangle, align=center, text width=\blockTextWidth cm, minimum width=\blockSize cm, minimum height=\blockSize cm] (box1) at (0, 0) {\hyperref[def:object]{\hl{\textbf{Object}}} \\ the system \\ the observer \\ seeks to understand};
\node[draw, rectangle, right of=box1, align=center, text width=\blockTextWidth cm, minimum width=\blockSize cm, minimum height=\blockSize cm] (box2) at (\blockDistance, 0) {\hyperref[def:explanation]{\hl{\textbf{Explanation}}} \\ \hyperref[def:interpretable]{\hl{interpretable}} \\ information enabling \\ the observer to refine \\ their mental model};
\node[draw, rectangle, right of=box2, align=center, text width=\blockTextWidth cm, minimum width=\blockSize cm, minimum height=\blockSize cm] (box3) at (2*\blockDistance, 0) {\textbf{Mental model} \\ of the observer \\ regarding the object};
% Arrows
\draw[-{Latex[length=3mm]}] (box1) -- node[midway, above, align=center, text width=\arrowTextWidth cm] {Information \\ extraction \\ through the use \\ of an \\ \hyperref[subsec:explainability_method]{\hl{\textbf{explainability}\\\textbf{method}}}} node[midway, below, align=center] {} (box2);
\draw[-{Latex[length=3mm]}] (box2) -- node[midway, above, align=center, text width=\arrowTextWidth cm] {\textbf{Interpretation} \\ by \\ the observer} node[midway, below, align=center] {} (box3);

% Legend: each entry is chained to the previous one via \legendGap /
% \legendIconTextGap, so the whole row reflows consistently if a label
% changes length or if the gaps are retuned.
\node[draw, rectangle, minimum width=0.3cm, minimum height=0.3cm] (legendIcon1) at (\legendXShift, \legendYShift) {};
\node[anchor=west, right=\legendIconTextGap of legendIcon1] (legendText1) {Information structures};

\node[inner sep=0pt, anchor=west, right=\legendGap of legendText1] (legendArrowStart) {};
\coordinate (legendArrowEnd) at ([xshift=0.4cm]legendArrowStart);
\draw[-{Latex[length=1.5mm]}] (legendArrowStart) -- (legendArrowEnd);
\node[anchor=west, right=\legendIconTextGap of legendArrowEnd] (legendText2) {Information transformation};

\node[anchor=west, right=\legendGap of legendText2] (legendIcon3) {\hl{~~}};
\node[anchor=west, right=\legendIconTextGap of legendIcon3] (legendText3) {Concepts discussed in more detail in this section};

% Brace: spans from just past box2 to just past box3, so it automatically
% follows \blockDistance and \blockSize instead of being retuned by hand.
\coordinate (C) at (\blockDistance+\blockSize/2+1, 2.6);
\coordinate (D) at (2*\blockDistance+\blockSize/2+1, 2.6);
\draw [decorate, decoration={brace, amplitude=10pt}, thick]
          (C) -- node[above=12pt, align=center] {Subjective analysis \\ by the \hyperref[def:observer]{\hl{observer}}} (D);
\end{tikzpicture}
    }
\caption{Explanation process in explainable artificial intelligence.}
\label{fig:explanation_process_illustration}
\end{figure}


Figure~\ref{fig:explanation_process_illustration} summarizes this whole explanation process, from the object of the explanation to the mental model refined by the observer, as a chain of three information structures, depicted as boxes and connected by arrows denoting a transformation of information:
\begin{enumerate}
    \item[(1)] The \textbf{object} of the explanation, that is, the system the observer seeks to understand.
    \item[(2)] The \textbf{explanation}, the interpretable information produced from the object.
    This transformation is performed through the use of an explainability method.
    \item[(3)] The \textbf{mental model} held by the observer regarding the object.
\end{enumerate}
A brace spanning boxes~(2) and~(3) indicates that this second half of the process, from the explanation to the resulting mental model, constitutes a subjective analysis carried out by the observer.

\subsection{Explainable Artificial Intelligence}

Now that the general framework of explanation has been established, we can examine how this concept applies in the context of \gls{ai}.
In the context of \gls{xai}, the object of explanation is an \gls{ai} system.
More precisely, \gls{xai} emerged as a response to the success of \gls{dl} models \cite{gunning2019darpa, arrieta2020explainable}.
In these architectures, information is represented in a distributed manner across the network through patterns of activation and connections between computational units.
Unlike symbolic approaches, individual components of \glspl{nn} cannot be directly associated with human-defined concepts.
From this starting point, we identify two types of approaches:
\begin{itemize}
    \item \textbf{Substitutive explainability:} In this approach, the underlying assumption is that there is nothing to explain within a \gls{nn}.
    There is no semantic meaning associated with each computational component, and therefore nothing to interpret.
    This corresponds to the so-called ``black box'' view of \glspl{nn}.
    Under this assumption, \gls{xai} consists of replacing as much of the \gls{dl} system as possible with alternative \gls{ai} methods \cite{rudin2019stop}.

    \item \textbf{Intrinsic explainability:} In this approach, it is argued that \glspl{nn} inherently contain mechanisms that can be explained.
    The limitation is considered to lie not in the networks themselves, but in the absence of adequate tools to analyze their functioning \cite{olah2020zoom}.
\end{itemize}

This manuscript considers methods from both approaches in order to provide a comprehensive overview of the \gls{xai} field.
According to the definition of explanation adopted in this manuscript (Definition~\ref{def:explanation}), we consider that any system can be explained.
As a consequence, the perspective adopted throughout this manuscript is that of intrinsic explainability.

\subsection{Explainability Method in XAI}
\label{subsec:explainability_method}

Now that the general framework of \gls{xai} has been established, we need to further detail the properties of an explanation in this context.
In our framework, an explanation in the field of \gls{xai} can be characterized by three components:
\begin{itemize}
    \item \textbf{Source:}\label{item:source} refers to the origin of the information used by the explainability method.
    The explainability method transforms the raw information produced by the \gls{ai} system into an explanation that can be understood by an observer.

    \item \textbf{Form:}\label{item:form} refers to the representation through which the explanation is communicated.
    The form defines how the information is presented to the observer and must allow human interpretation.
    An explanation can take many different forms, including text, images, computer code, and mathematical equations.

    \item \textbf{Subject:}\label{item:subject} refers to the aspect of the \gls{ai} system that is being explained.
    In \gls{xai}, this distinction is commonly associated with the difference between local and global explanations.
    A local explanation describes why an \gls{ai} system produced a specific output in a given situation.
    A global explanation aims to describe the general functioning of the \gls{ai} system.
\end{itemize}

These three properties of an explanation are entirely determined by the explainability method (Definition~\ref{def:explainability_method}) used to produce it.
In this context, an explainability method is a transformation process that maps a source of information into an explanation characterized by a specific form and subject.
Research in \gls{xai} focuses on developing new methods capable of performing this transformation in order to overcome the lack of interpretability associated with \gls{dl} models.